\documentclass[12pt,a4paper]{article}
\usepackage[spanish]{babel}
\usepackage[utf8]{inputenc}
\usepackage[T1]{fontenc}
\usepackage{geometry}
\usepackage{array}
\usepackage{longtable}
\usepackage{enumitem}
\usepackage{hyperref}
\usepackage{xcolor}
\geometry{margin=2.5cm}
\setlist[itemize]{leftmargin=*}
\setlist[enumerate]{leftmargin=*}
\title{\textbf{ServerFort}\\Infraestructura Ubuntu Server segura, monitorizada y automatizada}
\author{Hugo Poblaciones \and Sergi Blanca \and Unai Hervella\\[0.5em]\large 
\date{19 de octubre de 2026}

\begin{document}
\maketitle
\begin{center}
\large Proyecto de Grado Superior ASIX\\Memoria inicial (anteproyecto)
\end{center}
\newpage
\tableofcontents


\section{Introducción}
\subsection{De qué va el proyecto}
ServerFort consiste en montar una infraestructura de empresa pequeña con Ubuntu Server: servidores, red, seguridad y copias de seguridad. La idea es que, una vez montada, se mantenga casi sola y que nadie tenga que estar todo el día pendiente.

\subsection{El problema}
Muchas pymes ponen un servidor Linux sin un firewall bien configurado, sin backups automáticos y sin nadie que vigile si algo falla. Como normalmente no tienen un administrador de sistemas a jornada completa, los problemas se detectan tarde. Si entra alguien que no debe o se estropea un disco, ya se ha perdido tiempo o datos. Un fallo así puede parar el negocio varios días, por eso creemos que merece la pena resolverlo.

\subsection{Para quién está pensado}
\begin{itemize}
\item Una pyme o un departamento de IT pequeño que no tiene un administrador de sistemas a tiempo completo.
\item Una organización que necesite red, servicios (web, archivos, DNS\ldots) y acceso remoto seguro sin depender de terceros.
\item Nosotros mismos, ya que nos sirve para poner en práctica lo aprendido en el ciclo.
\end{itemize}

\subsection{Objetivos}
Queremos pasar de un Ubuntu Server sin configurar a una pequeña infraestructura de empresa. Para ello vamos a:
\begin{itemize}
\item Diseñar una red separada por zonas, con DMZ, servidores y acceso remoto.
\item Montar los servicios: DNS, DHCP, Apache, archivos y VPN.
\item Proteger el sistema con firewall, permisos y auditoría (e IDS/IPS si nos da tiempo).
\item Controlar el estado de las máquinas con Prometheus, logs y alertas.
\item Automatizar tareas con Ansible y scripts.
\item Hacer copias de seguridad y probar a restaurarlas, con un plan de recuperación ante desastres.
\end{itemize}

\subsection{Qué aporta de diferente}
No nos queremos quedar solo en instalar un servidor. Vamos a simular ataques y fallos (un escaneo con Nmap, un servicio caído, un disco lleno) para comprobar que la seguridad y las alertas funcionan de verdad. Además, todo se hará con Ansible para poder repetir la instalación desde cero en pocos minutos, y juntamos redes, seguridad, monitorización y recuperación en el mismo proyecto.

\subsection{Estado del proyecto}
Esta memoria se ha escrito antes de empezar el proyecto, así que todavía no hay resultados. Aquí explicamos qué queremos hacer, cómo lo vamos a organizar y cómo comprobaremos que funciona. Todo lo que aparece es una previsión y puede cambiar durante el curso.

Las decisiones concretas (programas que usaremos, direcciones IP, configuraciones) las iremos tomando en cada sprint. Los resultados, las pruebas y las conclusiones se añadirán en la memoria final.

\section{Alcance}
\begin{longtable}{|p{0.25\textwidth}|p{0.67\textwidth}|}
\hline
\textbf{Parte} & \textbf{Qué incluye}\\ \hline
Imprescindible (MVP) & Red separada, DNS/DHCP, servidor web, firewall, backups y monitorización básica con Prometheus y Grafana.\\ \hline
Ampliaciones & Suricata (IDS/IPS), LDAP/FreeIPA y servicios en contenedores Docker.\\ \hline
Fuera del alcance & Programar aplicaciones propias y montarlo en hardware real. Todo se hace con máquinas virtuales.\\ \hline
\end{longtable}

\subsection{Máquinas virtuales}
Usaremos dos máquinas virtuales:
\begin{itemize}
\item Ubuntu Server, donde están los servicios.
\item Ubuntu Client, desde donde hacemos todas las pruebas contra el servidor y nos conectamos a él de forma remota.
\end{itemize}

\section{Arquitectura}
La infraestructura se divide en estas partes:
\begin{itemize}
\item Internet: entrada y salida de tráfico.
\item Firewall: NAT, reglas, VPN y separación de redes.
\item DMZ: web y servicios que se publican.
\item LAN / usuarios: acceso a los recursos internos.
\item Servidores: DNS, LDAP, archivos, monitorización y backup.
\item Laboratorio: pruebas de seguridad controladas.
\end{itemize}
El esquema de red se hará en el Sprint 1 y se añadirá a la memoria final.

\section{Red (previsión)}
Vamos a dividir la red en cinco VLAN para separar cada función y que, si pasa algo en una zona, no afecte a todas.

\begin{longtable}{|p{0.12\textwidth}|p{0.25\textwidth}|p{0.34\textwidth}|p{0.2\textwidth}|}
\hline
\textbf{VLAN} & \textbf{Nombre} & \textbf{Para qué sirve} & \textbf{Red}\\ \hline
10 & Administración & Gestionar equipos y servidores & Por definir\\ \hline
20 & Usuarios & Acceso a recursos de la empresa & Por definir\\ \hline
30 & Servidores & Servicios internos & Por definir\\ \hline
40 & DMZ & Servicios expuestos & Por definir\\ \hline
50 & Laboratorio & Pruebas y pentesting & Por definir\\ \hline
\end{longtable}

Usaremos IPv4, DHCP, DNS y reglas de firewall entre las zonas. La idea es bloquear todo por defecto y abrir solo lo necesario.

Las subredes y las reglas concretas entre VLAN se definirán en el Sprint 1 y el Sprint 3.

\section{Servicios}
\begin{itemize}
\item DNS y DHCP: resolución de nombres y configuración automática de la red. Los programas concretos se elegirán en el Sprint 2.
\item SSH: administración remota con claves.
\item Web: Apache con HTTPS, en la DMZ. El dominio y el certificado se decidirán en el Sprint 2.
\item Archivos: Samba con permisos por usuario y grupo, para que cada usuario solo vea sus cosas.
\item VPN: WireGuard para conectarse desde fuera de forma segura.
\end{itemize}

\subsection{Configuración}
La configuración de cada servicio se documentará durante los Sprints 2 y 3 y se incluirá en la memoria final.

\section{Seguridad}
Medidas que vamos a aplicar:
\begin{itemize}
\item SSH solo con clave y sin acceso directo como root.
\item Instalar solo los servicios necesarios y mantener el sistema actualizado.
\item Usar sudo, grupos y permisos, con política de contraseñas.
\item Firewall con puertos mínimos y reglas por zonas.
\item Suricata para detectar tráfico sospechoso (ampliación).
\end{itemize}
La herramienta de firewall y su configuración se decidirán en el Sprint 3.

\section{Monitorización}
\begin{itemize}
\item Prometheus recoge métricas de los servidores y servicios.
\item Grafana muestra gráficas de CPU, RAM, disco, red y disponibilidad.
\item Alertmanager avisa cuando hay una incidencia.
\item Los logs se centralizan para poder revisarlos.
\end{itemize}
Los paneles y las alertas concretas se definirán en el Sprint 4.

\section{Copias de seguridad}
Seguiremos la regla 3-2-1: tres copias, en dos tipos de soporte y una de ellas fuera. Las copias serán diarias, semanales y mensuales, con un tiempo de retención. Se comprimen, se cifran y se guardan en un sitio distinto al servidor. También se comprobará automáticamente que las copias están bien, y probaremos a borrar datos del laboratorio y recuperarlos.

La herramienta, lo que se copia y el tiempo de retención se definirán en el Sprint 5.

\subsection{Plan de recuperación (Disaster Recovery)}
El plan de recuperación se redactará en el Sprint 5. Este plan consistirá en un plan por si una de las máquinas virtuales se rompe, recuperarla a través de una copia.

\section{Automatización}
\begin{itemize}
\item Ansible: usuarios, paquetes, SSH, firewall y servicios.
\item Scripts: backups, comprobaciones y mantenimiento.
\item Docker/Podman: para algunos servicios (ampliación).
\end{itemize}
Para montar todo en máquinas nuevas, se ejecutaría el playbook principal:
\begin{verbatim}
ansible-playbook -i inventory/hosts.ini site.yml
\end{verbatim}
La estructura del repositorio se definirá al empezar el proyecto.

\section{Pruebas de seguridad}
Haremos las pruebas siguiendo este ciclo: ataque o fallo $\rightarrow$ detección $\rightarrow$ corrección $\rightarrow$ validación (comprobar otra vez que está arreglado).
\begin{itemize}
\item Servidor caído.
\item Disco lleno.
\item Servicio parado.
\item Intentos de acceso por SSH.
\item Escaneo de puertos con Nmap.
\end{itemize}
Los resultados de las pruebas se añadirán en la memoria final.

\section{Organización}
\subsection{Reparto de tareas}
\begin{itemize}
\item Hugo Poblaciones: redes, firewall y VLAN.
\item Sergi Blanca: servicios (DNS, web, archivos) y Ansible.
\item Unai Hervella: seguridad, monitorización y backups.
\end{itemize}
Los roles pueden cambiar durante el curso según lo que necesitemos. De momento hemos pensado en la siguiente organización rotativa:

\begin{longtable}{|p{0.17\textwidth}|p{0.25\textwidth}|p{0.25\textwidth}|p{0.25\textwidth}|}
\hline
\textbf{Sprint} & \textbf{Coordinador} & \textbf{Resp. Técnico} & \textbf{Resp. Documentación}\\ \hline
1 & Hugo & Sergi & Unai\\ \hline
2 & Sergi & Unai & Hugo\\ \hline
3 & Unai & Hugo & Sergi\\ \hline
4 & Hugo & Sergi & Unai\\ \hline
5 & Sergi & Unai & Hugo\\ \hline
Final & Unai & Hugo & Sergi\\ \hline
\end{longtable}

Cabe recalcar que, según las dificultades que vayamos viendo, pueden variar los roles o irnos ayudando en las diferentes tareas.

\subsection{Sprints}
\begin{enumerate}
\item Diseño de la red, las VLAN y la arquitectura.
\item Servicios base: DNS, DHCP, web y archivos.
\item Firewall, VPN y hardening.
\item Monitorización, dashboards y alertas.
\item Backups y pruebas de ataques simulados.
\item Documentación, pruebas finales y presentación.
\end{enumerate}

\section{Problemas previstos}
\begin{itemize}
\item \textbf{Problemas técnicos:} Puede que algún servicio no arranque bien en la VM y nos lleve tiempo arreglarlo. Dejamos margen en cada sprint.
\item \textbf{Cosas nuevas:} Ansible y Suricata no los hemos usado antes. Empezaremos con tutoriales antes de aplicarlos.
\item \textbf{Falta de tiempo:} Si nos retrasamos, lo primero que quitaremos será el laboratorio de pentesting, para no afectar a lo imprescindible.
\end{itemize}
Los problemas reales que aparezcan se recogerán en la memoria final.

\section{Entrega y comprobación}
\subsection{Qué entregamos}
Las máquinas virtuales con la red y los servicios, el repositorio (Ansible, scripts, configuraciones y documentación), las pruebas de seguridad, los dashboards y alertas, los backups con su restauración, y esta memoria con la demo.

\subsection{Cómo demostramos que funciona}
\begin{enumerate}
\item Un usuario se autentica y solo accede a sus recursos.
\item Los servicios responden y están monitorizados.
\item El firewall bloquea el tráfico no autorizado.
\item Los intentos sospechosos generan registros y alertas.
\item Se hace un backup y se restaura.
\item Una incidencia simulada se detecta y se documenta.
\item La configuración se puede repetir con la automatización.
\end{enumerate}
Las capturas y evidencias de cada punto se recogerán durante el proyecto y se incluirán en la memoria final.

\section{Resultados esperados y conclusiones}
\subsection{Qué esperamos conseguir}
Al terminar el proyecto, esperamos tener una infraestructura funcionando que sea segura, esté vigilada, se pueda reconstruir con Ansible y permita recuperar los datos si algo falla. También esperamos aprender a defender, vigilar y recuperar un servidor, que es lo que se le pide a un administrador de sistemas en una empresa real.

\subsection{Conclusiones}
Las conclusiones se escribirán al terminar el proyecto, cuando tengamos resultados reales.

\end{document}
